\chapter{Literature Review}
\label{chap:literature-review}

% PLAN: the core of the January delivery, ~25--32 pages (W3--W8).
% Intro paragraph: positioning at the intersection of recourse, fairness and
% KANs, with causal modelling as the enabler (positioning figure). Explain that
% the chapter ends with a decision about the role of the KAN (decision model,
% causal model or recourse generator), taken up in Chapter~\ref{chap:proposed-approach}.
% Rule: every section ends with ``Summary and Implications'' tied to the thesis.

\section{Review Methodology}
\label{sec:lr-methodology}

% (W3) Protocol following \cite{kitchenham2007guidelines}; reporting following
% \cite{page2021prisma}; snowballing following \cite{wohlin2014guidelines}.

\subsection{Review Questions}
\label{subsec:lr-review-questions}

% The review's own questions (one per cluster), not the thesis RQs.

\subsection{Search Strategy}
\label{subsec:lr-search}

% Databases, preprint servers, venues, strings S1--S5 (table), search dates.

\subsection{Inclusion and Exclusion Criteria}
\label{subsec:lr-criteria}

% IC1--IC4, EC1--EC5, time windows.

\subsection{Study Selection}
\label{subsec:lr-selection}

% PRISMA flow diagram with the logged counts (finish in W8).

\section{Counterfactual Explanation and Recourse Generation Methods}
\label{sec:lr-recourse-methods}

% (W4) Taxonomy figure + comparison table:
% method | year | model access | causal? | actionability | plausibility | data | code.
% Seeds: \cite{karimi2022survey, verma2024counterfactual, guidotti2024counterfactual}.

\subsection{Optimisation-Based Methods}
\label{subsec:lr-optimisation}

% \cite{wachter2018counterfactual, ustun2019actionable, mothilal2020explaining}.

\subsection{Data-Manifold and Generative Methods}
\label{subsec:lr-generative}

% \cite{poyiadzi2020face, joshi2019towards, pawelczyk2020learning}.
% Note amortized generators (relevant to option O3).

\subsection{Causality-Aware Recourse}
\label{subsec:lr-causal-recourse}

% \cite{mahajan2019preserving, karimi2021algorithmic}.

\subsection{Robustness and Global Recourse}
\label{subsec:lr-robust-global}

% \cite{upadhyay2021towards, dominguezolmedo2022adversarial, rawal2020beyond};
% critical perspectives \cite{venkatasubramanian2020philosophical, barocas2020hidden}.

\subsection{Summary and Implications}
\label{subsec:lr-recourse-summary}

\section{Fairness of Algorithmic Recourse}
\label{sec:lr-fair-recourse}

% (W5) Comparison table:
% notion | definition | needs SCM? | mitigation proposed? | datasets.

\subsection{Group-Level Notions of Recourse Fairness}
\label{subsec:lr-group-recourse-fairness}

% Cost parity and burden \cite{gupta2019equalizing, sharma2020certifai};
% social burden \cite{barrainkua2026revisiting}; substantive equality
% \cite{bell2024fairness}.

\subsection{Individual and Causal Notions of Recourse Fairness}
\label{subsec:lr-causal-recourse-fairness}

% \cite{vonkugelgen2022fairness, ehyaei2023robustness}.

\subsection{Auditing and Mitigating Unfair Recourse}
\label{subsec:lr-recourse-mitigation}

% Subgroup audits \cite{kavouras2023fairness}; in-training vs. post-hoc
% mitigation (relevant to option O1).

\subsection{Summary and Implications}
\label{subsec:lr-fair-recourse-summary}

\section{Causal Modelling for Counterfactual Inference}
\label{sec:lr-causal-models}

% (W6) Comparison table:
% model | family | identifiability | mixed data? | interpretable mechanisms?

\subsection{Deep Structural Causal Models}
\label{subsec:lr-dscm}

% \cite{pawlowski2020deep, xia2023neural}.

\subsection{Flow-, Autoencoder- and Diffusion-Based Causal Models}
\label{subsec:lr-causal-generative}

% \cite{khemakhem2021causal, sanchezmartin2022vaca, javaloy2023causal, chao2023interventional}.

\subsection{Recourse under Imperfect Causal Knowledge}
\label{subsec:lr-imperfect-causal}

% \cite{karimi2020algorithmic}. Relevant to option O2: where does the graph come from?

\subsection{Summary and Implications}
\label{subsec:lr-causal-summary}

\section{Kolmogorov-Arnold Networks in Trustworthy Machine Learning}
\label{sec:lr-kan}

% (W7) Comparison table:
% variant/work | basis function | claim | tabular evaluation? | fairness/XAI/causal use?
% Surveys: \cite{somvanshi2024survey, hou2024comprehensive}.

\subsection{Variants and Efficiency}
\label{subsec:lr-kan-variants}

% \cite{li2024kolmogorov, yu2024kan} + variants found by the search.

\subsection{KANs on Tabular Data}
\label{subsec:lr-kan-tabular}

% \cite{poeta2024benchmarking, yadav2025kolmogorov}.

\subsection{Interpretability Claims and Evidence}
\label{subsec:lr-kan-interpretability}

% \cite{liu2025kan, liu2024kan}; contrast with \cite{rudin2019stop, agarwal2021neural}.

\subsection{KANs for Fairness, Explainability and Causality}
\label{subsec:lr-kan-trustworthy}

% \cite{priyadarshini2026learning} + works found by S5. State the gap
% (no KAN-based recourse work as of the search freeze), backed by the search log.

\subsection{Summary and Implications}
\label{subsec:lr-kan-summary}

\section{Datasets, Metrics and Benchmarking Practice}
\label{sec:lr-datasets}

% (W8) Datasets table:
% dataset | domain | n | features | sensitive attribute | actionable features | causal graph? | used by.
% \cite{lequy2022survey, becker1996adult, hofmann1994statlog, angwin2016machine,
%       wightman1998lsac, yeh2009comparisons}; benchmarking \cite{pawelczyk2021carla}.
% Also: which metrics and statistical tests are reported.

\section{Synthesis and Research Gaps}
\label{sec:lr-gaps}

% (W8) 1--2 pages. Cross-cluster synthesis -> explicit gaps -> final RQs
% (copy them back to Section~\ref{sec:intro-objectives}). End with the decision
% matrix for the KAN's role (options O1--O4).
